
\chapter*{Prólogo\markboth{Prólogo}{Prólogo}}

La idea de TierraViva nace de una situación cercana. En mi familia hay un pequeño huerto en el pueblo y, especialmente en épocas de calor, aparece siempre el mismo problema: no es posible revisarlo todas las semanas. A veces las plantas reciben menos agua de la necesaria y otras veces ocurre lo contrario, porque el riego depende de visitas puntuales, de la intuición y de la disponibilidad de cada momento.

La frase que terminó de dar forma a la idea surgió casi en broma. Un día, hablando sobre estos problemas, mi tío me dijo: \textit{``¿No vas a ser ingeniera? Pues arréglame el huerto''}. Aunque fue un comentario informal, planteaba un problema muy real: cómo conocer el estado de un cultivo cuando no se está físicamente allí y cómo actuar sobre el riego sin depender siempre de una revisión manual.

A partir de esa necesidad surgió la posibilidad de desarrollar un sistema capaz de medir el entorno, enviar información, visualizar el estado del cultivo y actuar sobre el riego de forma controlada. Lo que inicialmente parecía una solución sencilla fue tomando forma como un problema de ingeniería más completo, en el que intervienen sensorización, comunicaciones, sistemas embebidos, almacenamiento de datos, visualización y control.
